\clearpage
\phantomsection

\setcounter{chapter}{2}
\chapter[Cài đặt và thực nghiệm]{Cài đặt và thực nghiệm}

Chương này mô tả cách hiện thực hệ thống đã thiết kế ở Chương~2 trong Unity và Unity ML-Agents. Trọng tâm của chương không phải là liệt kê toàn bộ mã nguồn, mà là chỉ ra các thành phần quyết định đến khả năng tái lập thí nghiệm: môi trường phát triển, bộ sinh bản đồ, agent, cấu hình huấn luyện và phương pháp ablation.

\section{Môi trường phát triển}

Môi trường game được xây dựng bằng Unity 6 và TopDown Engine. Phần học tăng cường dùng Unity ML-Agents để kết nối scene Unity với trainer Python. Các thí nghiệm chính được chạy bằng bản build standalone ở chế độ không đồ họa, với nhiều environment song song để tăng tốc độ thu thập mẫu.

\begin{table}[H]
    \centering
    \small
    \caption{Môi trường triển khai và huấn luyện}
    \label{tab:development_environment}
    \begin{tabular}{|p{5cm}|p{8cm}|}
        \hline
        \textbf{Thành phần} & \textbf{Cấu hình sử dụng} \\
        \hline
        Game engine & Unity Editor 6000.1.3f1. \\
        \hline
        Gameplay framework & TopDown Engine v4.4, môi trường demo Koala2D được mở rộng cho bài toán RL. \\
        \hline
        Unity package & com.unity.ml-agents phiên bản 4.0.3. \\
        \hline
        Trainer & Unity ML-Agents Python trainer, thuật toán PPO, log qua TensorBoard. \\
        \hline
        Phần cứng & CPU 16 luồng, RAM 16GB, NVIDIA GeForce RTX 3050 6GB. \\
        \hline
        Chế độ chạy & Standalone build, --no-graphics, time-scale=20, thường dùng 8 environment song song; một số run dùng 16 environment. \\
        \hline
    \end{tabular}
\end{table}

Việc dùng build standalone thay vì chạy trực tiếp trong Editor giúp giảm nhiễu do giao diện và ổn định tốc độ mô phỏng. Với cấu hình 3 triệu bước, thời gian huấn luyện một run đầy đủ thường nằm ở thang nhiều giờ, trong đó phần tốn thời gian nhất là mô phỏng Unity và đồng bộ nhiều environment, không phải forward pass của mạng MLP.

\section{Hiện thực pipeline sinh bản đồ bằng PCGNN}

Pipeline bản đồ được tách thành hai giai đoạn. Giai đoạn offline chạy trong môi trường Python của PCGNN để sinh, đánh giá và chia bản đồ theo độ khó. Giai đoạn online chạy trong Unity thông qua TilemapGenerator để nạp bản đồ đã chọn vào episode huấn luyện. Cách tách này là cải tiến quan trọng so với phương án gọi generator trực tiếp từ Unity: quá trình train ổn định hơn, bản đồ có thể được kiểm tra trước, và người thiết kế vẫn có quyền chọn những map phù hợp nhất trước khi đưa vào tập huấn luyện chính.

Ở giai đoạn PCGNN, đề tài sử dụng genome đã huấn luyện \texttt{inctyseed0.pkl} và logic sinh từ notebook \texttt{improve-v2.ipynb}. Mỗi ô bản đồ được sinh theo thứ tự quét từ trái sang phải, từ trên xuống dưới; mạng nhận ngữ cảnh cục bộ $3 \times 3$ xung quanh ô hiện tại cùng bốn đầu vào ngẫu nhiên, sau đó quyết định ô đó là tường hay sàn. Kết quả sau đó được adapter chuyển sang text map của game.

\subsection{Các điều chỉnh so với PCGNN gốc}

Paper gốc của Beukman và cộng sự dùng PCGNN như một bộ sinh level độc lập và đánh giá chất lượng bằng các chỉ số như leniency, A* difficulty, compression distance và A* diversity~\cite{Beukman2022}. Trong khóa luận này, PCGNN được giữ lại ở vai trò bộ sinh ma trận nền, nhưng được mở rộng thành một pipeline phục vụ trực tiếp cho huấn luyện tác nhân trong Unity. Các điểm điều chỉnh chính được tóm tắt trong Bảng~\ref{tab:pcgnn_improvements}.

\begin{table}[H]
    \centering
    \small
    \caption{Các điều chỉnh của khóa luận so với cách dùng PCGNN gốc}
    \label{tab:pcgnn_improvements}
    \begin{tabular}{p{4.0cm} p{5.0cm} p{5.0cm}}
        \toprule
        \textbf{Thành phần} & \textbf{PCGNN gốc} & \textbf{Điều chỉnh trong khóa luận} \\
        \midrule
        Vai trò của generator & Sinh level để đánh giá chất lượng và độ đa dạng của nội dung. & Sinh tập bản đồ ứng viên offline để đưa vào curriculum learning trong Unity. \\
        Biểu diễn đầu ra & Ma trận tile của miền Maze/Mario trong thí nghiệm gốc. & Chuyển sang text map của game với ký hiệu \texttt{\#}, \texttt{.}, \texttt{P}, \texttt{E}. \\
        Xử lý vị trí người chơi và kẻ địch & Phụ thuộc vào biểu diễn của miền gốc. & Adapter đặt \texttt{P}/\texttt{E} trên cùng thành phần liên thông lớn nhất và chọn hai vị trí cách xa nhau. \\
        Kiểm tra hợp lệ & Đánh giá bằng các metric offline. & BFS xác nhận đường đi từ \texttt{P} tới \texttt{E}; TilemapGenerator kiểm tra lại vùng đi được khi nạp map. \\
        Chỉ số bổ sung & Paper báo cáo leniency, A* difficulty, compression distance và A* diversity. & Bổ sung \texttt{reachable\_ratio}, \texttt{dead\_end\_ratio}, \texttt{difficulty\_score} để phục vụ chọn map cho curriculum. \\
        Tổ chức dữ liệu & Không gắn với lesson của môi trường Unity. & Chia map thành Easy/Medium/Hard theo percentile để điều khiển độ khó huấn luyện. \\
        Tích hợp runtime & Không hướng tới gọi trực tiếp trong Unity. & Tách offline generator và online TilemapGenerator để giảm latency, tăng khả năng tái lập và cho phép kiểm tra map trước. \\
        \bottomrule
    \end{tabular}
\end{table}

Do đó, phần ``cải tiến'' trong khóa luận không nằm ở việc thay đổi hoàn toàn thuật toán NEAT của PCGNN, mà nằm ở lớp hậu xử lý, kiểm định, đo reachability và tích hợp bản đồ vào curriculum. Cách làm này phù hợp với mục tiêu của đề tài: tạo một tập bản đồ procedural đủ đa dạng nhưng vẫn kiểm soát được tính hợp lệ trước khi dùng để huấn luyện CombatAgent.

\subsection{Adapter và kiểm tra hợp lệ}

Adapter thực hiện ba việc. Thứ nhất, ma trận tile của PCGNN được chuyển sang ký tự \texttt{\#} và \texttt{.}. Thứ hai, hệ thống tìm thành phần liên thông lớn nhất của vùng sàn và đặt \texttt{P}/\texttt{E} ở hai ô cách xa nhau trong thành phần đó. Thứ ba, BFS được chạy lại để xác nhận đường đi từ \texttt{P} tới \texttt{E} tồn tại. Nhờ vậy, các map đưa vào Unity tránh được lỗi phổ biến của PCG như player bị nhốt, enemy nằm ở vùng không thể tiếp cận hoặc bản đồ không có đủ ô sàn hợp lệ.

Mỗi bản đồ hợp lệ được ghi ra file \texttt{.txt}, đồng thời một dòng chỉ số được lưu trong \texttt{metrics.csv}. Các chỉ số chính gồm:

\begin{itemize}
    \item \texttt{wall\_ratio}: tỷ lệ ô tường trên toàn bản đồ.
    \item \texttt{shortest\_path\_length}: độ dài đường đi ngắn nhất từ \texttt{P} tới \texttt{E}.
    \item \texttt{dead\_end\_ratio}: tỷ lệ ô đi được có số hàng xóm mở nhỏ hơn hoặc bằng một.
    \item \texttt{reachable\_ratio}: tỷ lệ vùng sàn thuộc thành phần liên thông lớn nhất.
    \item \texttt{astar\_difficulty}: mức độ A* phải mở rộng thêm node ngoài đường đi tối ưu.
    \item \texttt{difficulty\_score}: điểm tổng hợp dùng để xếp hạng tương đối.
\end{itemize}

Các đại lượng trên được dùng để đánh giá \emph{chất lượng hình học của ma trận bản đồ}, không phải hiệu năng cuối cùng của tác nhân. \texttt{wall\_ratio} kiểm soát mật độ vật cản; nếu quá thấp thì phòng gần như trống, nếu quá cao thì dễ tạo vùng bị chia cắt. \texttt{shortest\_path\_length} đo khoảng cách di chuyển tối thiểu giữa người chơi và kẻ địch, nên liên quan trực tiếp đến thời gian tiếp cận mục tiêu. \texttt{dead\_end\_ratio} phản ánh mức độ xuất hiện các nhánh cụt có thể khiến agent bị kẹt hoặc phải quay đầu. \texttt{reachable\_ratio} là chỉ số quan trọng nhất với bài toán combat, vì nó đo phần vùng sàn thực sự có thể khai thác để di chuyển và né tránh. \texttt{astar\_difficulty} mô tả độ phức tạp của cấu trúc đường đi theo thuật toán tìm đường, nhưng không đồng nghĩa hoàn toàn với độ khó chiến đấu. Cuối cùng, \texttt{difficulty\_score} là điểm xếp hạng tương đối, kết hợp các tín hiệu hình học để chia map thành Easy, Medium và Hard; điểm này chỉ có ý nghĩa so sánh trong cùng một batch sinh bản đồ.

\subsection{Chia tier cho curriculum}

Sau khi tính \texttt{difficulty\_score}, bản đồ được sắp xếp từ dễ đến khó. Trong phiên bản hiện tại, đề tài dùng cách chia theo percentile để tránh phụ thuộc vào ngưỡng thủ công quá cứng: 5\% bản đồ dễ nhất được đưa vào Easy, 5\% tiếp theo vào Medium và 90\% còn lại vào Hard. Cách chia này đặc biệt phù hợp khi đổi genome hoặc kích thước bản đồ, vì hệ thống luôn tạo được ba nhóm tương đối để người thiết kế xem lại.

Với cấu hình 14x14 dùng cho thử nghiệm tích hợp, pipeline đã sinh 1000 bản đồ và chia folder như Bảng~\ref{tab:pcgnn_map_generation_summary}.

\begin{table}[H]
    \centering
    \small
    \caption{Kết quả sinh và phân loại bản đồ PCGNN 14x14}
    \label{tab:pcgnn_map_generation_summary}
    \begin{tabular}{|p{5cm}|c|p{6.7cm}|}
        \hline
        \textbf{Chỉ số} & \textbf{Giá trị} & \textbf{Ý nghĩa} \\
        \hline
        Số bản đồ sinh ra & 1000 & Tập ứng viên để người thiết kế chọn map đưa sang Unity. \\
        \hline
        Số bản đồ duy nhất & 1000 & Không có trùng lặp theo nội dung file trong batch sinh. \\
        \hline
        Tỷ lệ đi tới hợp lệ & 1000/1000 & Adapter đặt \texttt{P}/\texttt{E} và BFS xác nhận map có đường đi. \\
        \hline
        Easy / Medium / Hard & 50 / 50 / 900 & Chia theo percentile 5\% / 5\% / 90\%. \\
        \hline
        Tỷ lệ tường trung bình & 0.322 & Mật độ vật cản ở mức vừa, phù hợp bản đồ 14x14. \\
        \hline
        Độ dài đường đi trung bình & 30.74 & Đường đi đủ dài để tạo bài toán di chuyển/tiếp cận. \\
        \hline
        Tỷ lệ ngõ cụt trung bình & 0.063 & Có một số cấu trúc bẫy/nhánh cụt nhưng không quá dày. \\
        \hline
        Điểm độ khó trung bình & 0.547 & Dùng để xếp hạng tương đối giữa các map trong cùng batch. \\
        \hline
    \end{tabular}
\end{table}

\subsection{Đối chiếu chỉ số PCGNN với paper gốc}

Để kiểm tra chất lượng bộ sinh bản đồ trước khi đưa sang Unity, đề tài chạy lại pipeline PCGNN trên năm seed độc lập (seed 0--4, mỗi seed sinh 100 bản đồ) và tính các chỉ số thường dùng trong paper gốc của Beukman và cộng sự~\cite{Beukman2022}. Bảng~\ref{tab:pcgnn_four_seed_metrics} trình bày kết quả theo từng seed. Trong đó, \textbf{seed 0} chính là genome production \texttt{inctyseed0.pkl} được dùng để sinh bộ map huấn luyện chính cho Unity. Chỉ số \texttt{reachable\_ratio} được bổ sung trong khóa luận để đo tỷ lệ vùng đi được thuộc thành phần liên thông lớn nhất; chỉ số này không được báo cáo trực tiếp trong paper gốc, nhưng hữu ích khi tuyển chọn map cho môi trường game.

\begin{table}[H]
    \centering
    \scriptsize
    \setlength{\tabcolsep}{4pt}
    \renewcommand{\arraystretch}{1.15}
    \caption{Kết quả PCGNN trên năm seed độc lập (100 bản đồ mỗi seed)}
    \label{tab:pcgnn_four_seed_metrics}
    \begin{tabular}{c c c c c c c}
        \toprule
        \textbf{Seed} & \textbf{Solv.} & \textbf{Reach.} & \textbf{Leniency} & \textbf{A* Diff.} & \textbf{Comp. Dist.} & \textbf{A* Div.} \\
        \midrule
        0 (production) & 79.0\% & 0.9216 & 0.5802 & 0.3987 & 0.4945 & 0.2051 \\
        1 & 92.0\% & 0.8720 & 0.6622 & 0.3138 & 0.4899 & 0.1240 \\
        2 & 78.0\% & 0.9480 & 0.5876 & 0.3026 & 0.4837 & 0.1993 \\
        3 & 74.0\% & 0.9740 & 0.6902 & 0.3582 & 0.4550 & 0.2180 \\
        4 & 94.0\% & 0.8940 & 0.8722 & 0.3244 & 0.4063 & 0.1014 \\
        \midrule
        Trung bình $\pm$ độ lệch chuẩn & 83.4\% $\pm$ 8.0\% & 0.9220 $\pm$ 0.047 & 0.6785 $\pm$ 0.1057 & 0.3396 $\pm$ 0.0349 & 0.4659 $\pm$ 0.0328 & 0.1696 $\pm$ 0.0474 \\
        \bottomrule
    \end{tabular}
\end{table}

Bảng~\ref{tab:pcgnn_paper_comparison} so sánh các chỉ số trùng với paper gốc trên miền Maze. Kết quả cho thấy độ dễ của bản đồ và A* diversity nằm gần vùng báo cáo của PCGNN gốc, trong khi A* difficulty cao hơn đáng kể. Điều này phù hợp với mục tiêu của bản cải tiến trong khóa luận: ưu tiên tạo bản đồ có cấu trúc liên thông và đủ đa dạng để phục vụ curriculum, sau đó dùng adapter/BFS để đảm bảo bản đồ đưa vào Unity có đường đi hợp lệ.

\begin{table}[H]
    \centering
    \small
    \caption{So sánh kết quả PCGNN với paper gốc trên miền Maze (trung bình $\pm$ độ lệch chuẩn trên 5 seeds)}
    \label{tab:pcgnn_paper_comparison}
    \begin{tabular}{l c c p{5.8cm}}
        \toprule
        \textbf{Chỉ số} & \textbf{Kết quả khóa luận} & \textbf{PCGNN gốc} & \textbf{Nhận xét} \\
        \midrule
        Leniency & 0.6785 $\pm$ 0.1057 & 0.70 $\pm$ 0.08 & Gần với paper gốc. \\
        A* Difficulty & 0.3396 $\pm$ 0.0349 & 0.06 $\pm$ 0.08 & Cao hơn paper, cho thấy đường đi và cấu trúc map phức tạp hơn. \\
        Compression Distance & 0.4659 $\pm$ 0.0328 & 0.488 $\pm$ 0.002 & Thấp hơn nhẹ về trung bình và dao động lớn hơn. \\
        A* Diversity & 0.1696 $\pm$ 0.0474 & 0.13 $\pm$ 0.17 & Nằm trong vùng kết quả của paper gốc. \\
        Reachability & 0.9220 $\pm$ 0.047 & -- & Chỉ số bổ sung của khóa luận sau bước hậu xử lý adapter/BFS; mục tiêu tuyển chọn map cho Unity là khoảng 0.90 $\pm$ 0.10. \\
        \bottomrule
    \end{tabular}
\end{table}

Một điểm cần làm rõ là quy ước spawn và chỉ số solvability của paper gốc \textbf{không phải mục tiêu cần đạt} cho khóa luận này, do hai bài toán có tính chất khác nhau. Paper gốc đặt PCGNN trong miền Maze: bài toán là tìm đường từ một điểm xuất phát \emph{cố định} ở góc trên trái tới một điểm đích \emph{cố định} ở góc dưới phải, nên solvability 100\% là điều kiện đủ và là chỉ số đánh giá tự nhiên. Trong các kết quả tái lập baseline, nhiều bản đồ thỏa điều kiện này bằng cấu trúc đường đi kiểu chữ L hoặc corridor bám theo biên (Hình~\ref{fig:pcgnn_l_pattern_reference}): có đường nối hai góc nhưng phần còn lại của bản đồ có thể là vùng tường đặc hoặc các ngõ cụt rời rạc.

Bài toán của khóa luận là một game hành động combat trong phòng sinh theo thủ tục, không phải maze pathfinding. Tác nhân không di chuyển từ một góc cố định đến một góc cố định, mà phải \emph{di chuyển linh hoạt, né đạn, tiếp cận và tiêu diệt nhiều kẻ địch} trên toàn bộ vùng sàn. Do đó, ba hệ quả sau làm cho quy ước của paper gốc không còn phù hợp:

\begin{itemize}
    \item \textbf{Vị trí spawn không cố định ở góc.} Người chơi và kẻ địch có thể được đặt ở nhiều ô sàn khác nhau theo lesson (fixed spawn ở các ô sàn trung tâm hoặc random spawn trong toàn vùng đi được). Vì vậy, kiểm tra đường đi giữa hai góc cụ thể của bản đồ không phản ánh đúng yêu cầu thực tế.
    \item \textbf{Solvability 100\% không phải mục tiêu.} Đề tài không yêu cầu mọi bản đồ đầu ra của PCGNN đều hợp lệ -- bản đồ không đạt sẽ bị loại ở bước hậu xử lý. Cái cần kiểm tra là \emph{vùng đi được liên thông đủ rộng} để chứa nhiều cặp \texttt{P}/\texttt{E} hợp lệ và đủ không gian chiến đấu.
    \item \textbf{Cần chỉ số đo chất lượng vùng sàn, không phải đường đi.} Vì vậy khóa luận thay solvability bằng \texttt{reachable\_ratio} (tỷ lệ vùng sàn thuộc thành phần liên thông lớn nhất) -- đây là chỉ số phù hợp hơn với combat game vì nó đo khả năng \emph{di chuyển tự do} chứ không chỉ đo \emph{có đường đi không}.
\end{itemize}

\begin{figure}[H]
    \centering
    \includegraphics[width=0.95\textwidth]{pcgnn_l_pattern_reference.png}
    \caption{Ví dụ các bản đồ PCGNN kiểu Maze có đường đi hợp lệ từ góc trên trái tới góc dưới phải nhưng dễ hình thành L-pattern/corner-path}
    \label{fig:pcgnn_l_pattern_reference}
\end{figure}

Vì vậy, mục tiêu của khóa luận không phải chỉ là sinh ra một ma trận có đường đi hợp lệ theo định nghĩa Maze, mà là tạo các bản đồ có thể được đưa vào Unity để agent thực sự chơi được: có vùng sàn đủ liên thông, vị trí \texttt{P}/\texttt{E} hợp lệ, có không gian di chuyển và có độ khó phù hợp cho curriculum. Hình~\ref{fig:pcgnn_inctyseed0_maps} minh họa một nhóm bản đồ được sinh từ genome \texttt{inctyseed0.pkl} sau bước chuyển sang định dạng của game. Các map này vẫn giữ ký hiệu \texttt{P} và \texttt{E}, nhưng được đánh giá thêm bằng BFS, reachability và difficulty score trước khi đưa vào các folder Easy/Medium/Hard.

\begin{figure}[H]
    \centering
    \includegraphics[width=0.95\textwidth]{pcgnn_inctyseed0_maps.png}
    \caption{Một số bản đồ 14x14 sinh từ \texttt{inctyseed0.pkl} và được chuyển sang định dạng \texttt{\#}, \texttt{.}, \texttt{P}, \texttt{E} của game}
    \label{fig:pcgnn_inctyseed0_maps}
\end{figure}

\subsection{Nạp bản đồ vào Unity}

Tập tin TilemapGenerator.cs chịu trách nhiệm chuyển bản đồ văn bản thành tilemap, chọn bản đồ theo bài học, chọn vị trí sinh và sinh kẻ địch. Các bản đồ sau khi tuyển chọn được đặt trong ba thư mục Easy, Medium và Hard. Các bài học trong curriculum không sinh bản đồ mới từ đầu, mà lấy mẫu từ nhóm bản đồ tương ứng để vừa có đa dạng, vừa giữ khả năng tái lập.

Một phần logic chọn map theo curriculum có thể tóm tắt như sau:

\begin{lstlisting}
int diff = Mathf.RoundToInt(
    Academy.Instance.EnvironmentParameters
        .GetWithDefault("difficulty", defaultDifficulty));

switch (Mathf.Clamp(diff, 0, 3))
{
    case 0:
        pool = easyMaps;   // Easy + Fixed spawn
        break;
    case 1:
        pool = mediumMaps; // Medium + Fixed spawn
        break;
    case 2:
        pool = hardMaps;   // Hard + Fixed spawn
        break;
    default:
        pool = hardMaps;   // Hard + Random spawn
        break;
}
\end{lstlisting}

Sau khi chọn bản đồ, TilemapGenerator phân tích lưới ký tự, xác định ô sàn và ô tường, rồi dùng BFS để kiểm tra vùng có thể đi tới từ vị trí sinh của người chơi. Với chế độ sinh cố định, hệ thống ưu tiên ô sàn có số lượng ô có thể đi tới lớn để tránh các episode lỗi như người chơi hoặc kẻ địch bị cô lập. Với chế độ sinh ngẫu nhiên, hệ thống chọn từ các ô sàn hợp lệ và có thể giới hạn số kẻ địch được sinh để giảm phương sai quá lớn ở đầu huấn luyện.

\subsection{Các cải tiến đã thực hiện}

So với phiên bản chỉ dùng map văn bản được chuẩn bị thủ công, pipeline hiện tại có các cải tiến chính sau:

\begin{itemize}
    \item Bổ sung nguồn sinh map bằng PCGNN, giúp tạo nhanh số lượng lớn bản đồ ứng viên thay vì phụ thuộc hoàn toàn vào thiết kế thủ công.
    \item Thêm adapter chuyển ma trận PCGNN sang format \texttt{\#}, \texttt{.}, \texttt{P}, \texttt{E} của game.
    \item Tự động đặt \texttt{P}/\texttt{E} trên cùng thành phần liên thông để giảm lỗi bản đồ không thể hoàn thành.
    \item Thêm bộ đánh giá độ khó bằng BFS, A*, tỷ lệ tường, độ dài đường đi và tỷ lệ ngõ cụt.
    \item Thêm cơ chế chia tier theo percentile, tạo folder Easy, Medium và Hard để Unity dùng trực tiếp trong curriculum.
    \item Giữ pipeline offline để người thiết kế có thể duyệt map trước khi đưa vào huấn luyện chính.
\end{itemize}

\section{Hiện thực tác nhân CombatAgent}

Tập tin CombatAgent.cs là thành phần trung tâm của hệ thống. Tác nhân kế thừa lớp Agent của ML-Agents và cài đặt lại ba nhóm hàm quan trọng: CollectObservations, WriteDiscreteActionMask và OnActionReceived.

\subsection{Thu thập quan sát}

CollectObservations tạo vector 207 chiều đã mô tả ở Bảng~\ref{tab:obs_207}. Hàm này truy vấn trạng thái nhân vật, vũ khí, kẻ địch, projectile và các tia dò tường. Mỗi nhóm có số chiều cố định, vì vậy các đối tượng không tồn tại được thay bằng 0.

\begin{lstlisting}
// Pseudocode extracted from CombatAgent.cs
Collect player health, ammo, cooldown, velocity
Collect global visible-enemy and visible-bullet ratios
For up to 3 nearest enemies:
    add relative position, velocity, distance, health, threat
For up to 10 nearby bullets:
    add relative position, velocity, distance, time-to-impact
For 16 wall rays:
    add normalized hit distance
\end{lstlisting}

Thiết kế quan sát này làm cho chính sách nhìn thấy những yếu tố trực tiếp quyết định hành vi: khoảng cách tới địch, mức nguy hiểm của đạn, cooldown vũ khí và tường xung quanh. Đây là mức trừu tượng phù hợp với mục tiêu nghiên cứu curriculum và exploration, vì agent vẫn phải học chiến thuật nhưng không phải học xử lý ảnh.

\subsection{Xử lý hành động và action mask}

OnActionReceived đọc bốn nhánh hành động: moveX, moveY, combat và aim. Hai nhánh đầu được giải mã thành vector di chuyển, nhánh aim thành hướng ngắm, còn nhánh combat quyết định không hành động, tấn công hoặc dash. Hệ thống có mặt nạ hành động cho các hành động không hợp lệ, ví dụ tắt hành động tấn công khi vũ khí chưa sẵn sàng hoặc tắt dash khi dash đang hồi.

Cách tách nhánh này có hai lợi ích. Thứ nhất, nó cho phép agent đồng thời di chuyển và ngắm, gần với điều khiển game hành động thật. Thứ hai, action mask giảm lượng mẫu bị lãng phí bởi các hành động chắc chắn không thực thi được trong frame hiện tại.

\subsection{Ghi log chỉ số tùy biến}

Ngoài các scalar mặc định của ML-Agents như tổng phần thưởng, entropy, value loss và policy loss, CombatAgent ghi thêm nhiều chỉ số theo episode:

\begin{itemize}
    \item Số kẻ địch bị tiêu diệt trung bình mỗi episode.
    \item Tỷ lệ dọn sạch phòng.
    \item Tổng sát thương gây ra.
    \item Tổng sát thương nhận vào.
    \item Tỷ lệ tấn công đúng hướng.
    \item Tỷ lệ dash hữu ích.
\end{itemize}

Các chỉ số này quan trọng vì reward tổng hợp không đủ để giải thích hành vi. \textit{Cumulative reward} đo hiệu quả tổng hợp của chính sách theo đúng hàm thưởng đã thiết kế, nhưng có thể tăng nhờ nhiều nguyên nhân khác nhau. Vì vậy, \textit{clear rate} được dùng để đo trực tiếp tỷ lệ episode mà agent hoàn thành phòng; số kẻ địch bị tiêu diệt và sát thương gây ra phản ánh năng lực tấn công; sát thương nhận vào phản ánh mức độ né tránh và sống sót; tỷ lệ tấn công đúng hướng cho biết agent có học được quan hệ giữa vị trí địch và hướng ngắm hay không; tỷ lệ dash hữu ích cho biết agent có dùng dash như một hành động chiến thuật thay vì kích hoạt ngẫu nhiên hay không. Entropy của policy được đọc kèm các chỉ số này để đánh giá mức độ ổn định của chính sách: entropy cao nhưng reward và clear rate thấp thường cho thấy agent còn hành động quá ngẫu nhiên, trong khi entropy giảm cùng với reward/clear rate tăng là dấu hiệu chính sách đang hội tụ.

\section{Cấu hình huấn luyện PPO}

Cấu hình PPO baseline được chuẩn hóa quanh run46 để làm điểm so sánh cho các ablation curriculum. So với bản thử nghiệm cũ 200 nghìn bước, các run chính tăng lên 1--3 triệu bước vì môi trường 207 chiều và bản đồ procedural khó hơn đáng kể.

\begin{table}[H]
    \centering
    \small
    \caption{Cấu hình PPO baseline dùng trong các run chính}
    \label{tab:ppo_baseline_config}
    \begin{tabular}{|p{4cm}|c|p{7cm}|}
        \hline
        \textbf{Tham số} & \textbf{Giá trị} & \textbf{Ghi chú} \\
        \hline
        trainer\_type & ppo & Thuật toán chính. \\
        \hline
        batch\_size & 2048 & Batch lớn để ổn định gradient trong môi trường nhiều phương sai. \\
        \hline
        buffer\_size & 20480 & Tỷ lệ buffer:batch là 10:1. \\
        \hline
        learning\_rate & $3.0 \times 10^{-5}$ & Linear schedule. \\
        \hline
        beta & 0.004 & Hệ số entropy bonus. \\
        \hline
        epsilon & 0.1 & PPO clipping. \\
        \hline
        lambd & 0.95 & GAE. \\
        \hline
        num\_epoch & 2 & Số epoch trên mỗi buffer. \\
        \hline
        hidden\_units & 256 & MLP actor-critic. \\
        \hline
        num\_layers & 2 & Hai lớp ẩn. \\
        \hline
        time\_horizon & 64 & Độ dài trajectory trước khi update. \\
        \hline
        max\_steps & 1M hoặc 3M & 1M cho baseline nhanh, 3M cho curriculum chính. \\
        \hline
    \end{tabular}
\end{table}

Các cấu hình phần thưởng nội tại bổ sung tín hiệu curiosity (ICM) hoặc RND. Các cấu hình LSTM thêm bộ nhớ có kích thước 256 và chuỗi quan sát dài 64 bước. Các cấu hình curriculum thêm tham số environment\_parameters.difficulty để trainer tự chuyển bài học dựa trên phần thưởng đã làm mượt.

\section{Thiết kế ablation}

Ablation được thiết kế theo hướng thay đổi một yếu tố chính ở mỗi nhóm run. Mục tiêu không phải tìm một tổ hợp tối ưu tuyệt đối, mà là xác định yếu tố nào thực sự tạo cải thiện trên môi trường procedural combat.

\begin{table}[H]
    \centering
    \small
    \caption{Các run ablation chính}
    \label{tab:ablation_matrix_design}
    \begin{tabular}{|p{2cm}|p{5.1cm}|c|p{5.2cm}|}
        \hline
        \textbf{Run} & \textbf{Thiết lập} & \textbf{Steps log} & \textbf{Mục đích} \\
        \hline
        run34 & PPO baseline & 1.0M & Mốc so sánh không curriculum, không intrinsic reward. \\
        \hline
        run35 & PPO + ICM & 1.0M & Kiểm tra intrinsic motivation theo Pathak et al. \cite{Pathak2017}. \\
        \hline
        run36 & PPO + RND & 1.0M & Kiểm tra intrinsic motivation theo Burda et al. \cite{Burda2018}. \\
        \hline
        run42 & PPO + LSTM & 1.0M & Kiểm tra lợi ích của memory architecture \cite{Hochreiter1997}. \\
        \hline
        run46 & PPO + Curriculum & 3.0M & Thí nghiệm curriculum chính. \\
        \hline
        run47 & Curriculum + LSTM & 3.0M & Kiểm tra LSTM có cộng hưởng với curriculum hay không. \\
        \hline
        run50 & Curriculum + ICM & 3.0M & Kiểm tra ICM có phục hồi khi được scaffold bằng curriculum hay không. \\
        \hline
        run54 & 4-tier fixed-to-random & 3.0M & Kiểm tra giảm phương sai spawn đầu training. \\
        \hline
        run55 & Curriculum + reward v2 & 3.0M & Kiểm tra thay đổi thang reward và tăng trọng số outcome. \\
        \hline
    \end{tabular}
\end{table}

Do giới hạn thời gian và tài nguyên, mỗi run trong ma trận trên được chạy một seed. Các run curriculum chính gồm run46, run47, run50, run54 và run55 đều có log đến 3.0M bước, nên Chương~4 có thể so sánh trực tiếp trên cùng thang huấn luyện. Dù vậy, các kết quả vẫn được diễn giải theo hướng thực nghiệm định hướng: so sánh cẩn trọng các xu hướng lớn, tránh khẳng định quá mức với các chênh lệch nhỏ khi chưa có nhiều seed.

\section{Trích xuất dữ liệu}

Sau khi huấn luyện, Unity ML-Agents ghi kết quả của từng run vào các event file của TensorBoard trong thư mục Assets/results. Các event file này không thuận tiện để đưa trực tiếp vào bảng và biểu đồ, vì vậy khóa luận sử dụng script scripts/extract\_run\_metrics.py để đọc log và ghi số liệu ra các file CSV.

Script sử dụng EventAccumulator của TensorBoard để đọc các scalar như cumulative reward, policy entropy, cleared-room rate, số kẻ địch tiêu diệt, sát thương gây ra và các chỉ số curiosity. Kết quả trích xuất được ghi ra hai file dữ liệu:

\begin{itemize}
    \item data/run\_metrics.csv: lưu các chỉ số tổng hợp của từng run.
    \item data/rolling\_reward\_200k.csv: lưu reward trung bình trượt với cửa sổ 200 nghìn bước tại các mốc huấn luyện.
\end{itemize}

Chương~3 chỉ mô tả quy trình trích xuất để đảm bảo khả năng tái lập. Danh sách chỉ số, ý nghĩa của từng cột và các giá trị số cụ thể được trình bày ở Chương~4, nơi các kết quả được phân tích trực tiếp.
